\documentclass[sigconf,anonymous,review,natbib=false]{acmart}

% Supplemental material -- reproducibility appendix for the anonymous submission.
\settopmatter{printacmref=false}
\renewcommand\footnotetextcopyrightpermission[1]{}
\pagestyle{plain}

\acmConference[KDD AgenticAI Eval 2026]{Workshop on Evaluation and Trustworthiness of Agentic AI at KDD 2026}{August 2026}{Jeju, Republic of Korea}
\acmYear{2026}
\copyrightyear{2026}
\acmISBN{}
\acmDOI{}

\usepackage{booktabs}
\usepackage{listings}
\usepackage[most]{tcolorbox}
\usepackage{enumitem}
\usepackage{xcolor}
\lstset{breaklines=true, basicstyle=\small\ttfamily, language=Python}
\setlength{\emergencystretch}{3em}
\sloppy

\begin{document}

\title{Supplemental Material:\\Beyond Task Completion: Auditing Evolving Tool Libraries in Agentic AI}
\author{Anonymous Author(s)}
\affiliation{%
  \institution{Submission \#TBD}
  \country{}
}
\maketitle

This supplement provides reproducibility material for the strict
verified-subset pilot reported in the main paper. All tables and
manifest files are regenerated from \texttt{results\_canonical/} by the
scripts described in Section~4 of the main paper
(\texttt{build\_canonical.py}, \texttt{stats\_report.py},
\texttt{make\_tables.py}, \texttt{make\_figures.py}).

\section{Reproducibility Summary}
\label{app:reproducibility}

\paragraph{Model and decoding.}
All strict-pilot runs use Claude Haiku~4.5
(snapshot \texttt{claude-haiku-4-5-20251001}), accessed via the
Anthropic API with temperature 1.0 and the model default top-$p$.
Tool execution is sandboxed in-process with a five-second wall-clock
limit; harness state resets between sessions so library accumulation
is per-session.

\paragraph{Seeds and sessions.}
Three integer seeds (\texttt{0, 1, 2}) $\times$ eight verified-subset
sessions $\times$ five protocols yields 120 session rows. Each session
contributes a paired observation per protocol in the contrast
analyses.

\paragraph{Statistics.}
Per-session paired differences are bootstrap-resampled $B{=}10{,}000$
times with seed \texttt{20260604}; 95\% confidence intervals are
percentile CIs over the bootstrap distribution; the three
pre-specified contrasts (Table~3 in the main paper) are
Benjamini--Hochberg adjusted.

\paragraph{Canonical artifacts.}
The release includes the canonical run manifest
(\texttt{run\_manifest.jsonl}, 120 rows), tool manifest
(\texttt{tool\_manifest.jsonl}, system-level summary), aggregate
(\texttt{aggregate.json}), audit report
(\texttt{audit\_report.json}), statistical report
(\texttt{statistical\_report.json}), and the pre-specified contrasts
(\texttt{claims.json}). Every cell in the main-paper tables can be
derived from these files by the scripts above.

\section{Audit Coverage Detail}
\label{app:audit}

Table~\ref{tab:audit-detail} shows the per-session decomposition of
the strict subset. Across the eight verified-subset sessions, 51
graded task decisions per system pass survive into the strict TC
denominator. Seed tasks are explicitly excluded from the denominator
by design: they exist to expose the agent to the seed library before
any creation and carry no graded outcome.

\begin{table}[h]
\centering\footnotesize
\caption{Per-session strict coverage. Columns are role-level verified/total.
Seed counts are reported but excluded from TC (no graded outcome).}
\label{tab:audit-detail}
\begin{tabular}{lcccccc}
\toprule
\textbf{Session} & \textbf{Seed} & \textbf{Gap} & \textbf{Var.} & \textbf{Comp.} & \textbf{Reg.} & \textbf{Adv.} \\
\midrule
data\_transform\_s1 & 0/3 & 2/2 & 2/2 & 0/1 & 1/1 & 1/2 \\
data\_transform\_s2 & 0/3 & 2/2 & 2/2 & 1/1 & 1/1 & 1/2 \\
data\_transform\_s3 & 0/3 & 2/2 & 2/2 & 1/1 & 1/1 & 1/2 \\
data\_transform\_s4 & 0/3 & 2/2 & 2/2 & 1/1 & 1/1 & 2/2 \\
data\_transform\_s5 & 0/3 & 2/2 & 2/2 & 1/1 & 1/1 & 2/2 \\
numerical\_s1       & 0/3 & 2/2 & 2/2 & 0/1 & 1/1 & 1/2 \\
numerical\_s2       & 0/3 & 1/2 & 1/2 & 0/1 & 1/1 & 1/2 \\
numerical\_s3       & 0/3 & 2/2 & 2/2 & 0/1 & 1/1 & 0/2 \\
\midrule
Strict subset       & 0/24 & 15/16 & 15/16 & 4/8 & 8/8 & 9/16 \\
\bottomrule
\end{tabular}
\end{table}

\section{Baseline Fairness Ledger}
\label{app:fairness}

Three of the five protocols evaluated in the strict pilot are
adaptations of prior published systems to a shared open-task harness;
No-Evolution and One-Shot are native to this harness.
Table~\ref{tab:baseline_fairness} lists what was preserved from each
original design and where the adaptation deviates. The
\texttt{-style} suffix on EvoSkill-style, ToolMaker-style, and
CREATOR-style is deliberate: no underperformance in this pilot should
be read as a verdict on the original methods.

\begin{table}[h]
\centering\footnotesize
\renewcommand{\arraystretch}{1.2}
\begin{tabular}{p{0.10\textwidth}p{0.25\textwidth}p{0.25\textwidth}p{0.25\textwidth}}
\toprule
\textbf{Protocol} & \textbf{Native setting} & \textbf{Preserved faithfully} & \textbf{Adapted / fairness note} \\
\midrule
EvoSkill-style & Skill-library agent with NL strategy proposer/prompt generator & Skill-proposer and prompt-generator templates; strategy-folder structure & Strategies injected as prompt hints, not executed. Characterises strategy-level evolution, not native EvoSkill performance. \\
CREATOR-style & Decision/execution split tool creation, with native reference tests & Two-phase decide-then-create-and-use control flow; per-tool function synthesis & Our open tasks supply no native reference tests; identical verifier set for all systems. \\
ToolMaker-style & Repo-grounded with reference tests; diagnose-and-rewrite loop gated by provided tests & Two-stage diagnose-then-rewrite prompt structure; sandbox execution of auto-tests & Open tasks supply no reference tests to gate the loop; tool creation rate is therefore conservative and reflects the harness regime, not the model. \\
\bottomrule
\end{tabular}
\caption{Baseline fairness ledger for adapted protocols. All systems
share the identical task stream, seed library, model
(Claude Haiku~4.5), decoding parameters, and verifier set; only the
tool-creation control flow varies.}
\label{tab:baseline_fairness}
\end{table}

\section{Worked Trace Example: a Self-Validating Tool that Fails Hidden Tests}
\label{app:trace}

The main paper points to a worked case in \texttt{data\_transform\_s1}
where the audit layer flags a tool that TC scores as successful. This
appendix expands the case.

The shared capability \texttt{abr\_binary\_decode} appears in five
session-1 tasks: \texttt{gap\_1} (which invites creation of the
capability), \texttt{variant\_1}, \texttt{compose\_1},
\texttt{regress\_1}, and \texttt{adversarial\_1} (each of which can
reuse the tool created at \texttt{gap\_1}). In the preserved-source
Sonnet pilot, a single tool (\texttt{decode\_abr\_format}) was
synthesised alongside an agent-written test file. The tool passed
its self-tests at $C{=}0.50$; the session-level TC reached $0.95$.

When the same tool source is replayed against the held-out conformance
suite tied to the capability (six unit inputs and six adversarial
inputs, none seen by the agent), correctness is $0.00$ and
robustness/generality are $0.50/0.00$. The tool is implementing a
parser that structurally agrees with its own training-style examples
but does not implement the capability's contract.

\begin{tcolorbox}[colframe=black!70, colback=gray!5, boxrule=0.5pt]
\textbf{What the audit layer surfaces.}
\begin{itemize}[noitemsep,leftmargin=1.2em]
  \item Session TC: $0.95$ -- a near-perfect session.
  \item In-session reuse: 1 correct, 1 incorrect -- ambiguous on its own.
  \item Hidden-test correctness on the same artifact: $0.00$ -- the tool fails every held-out input for the capability it claims to implement.
\end{itemize}
\end{tcolorbox}

The header pattern (TC high, in-session reuse mixed, hidden-test
correctness near zero on the underlying artifact) is exactly the
failure mode the framework is designed to expose. Per-tool conformance
testing under the strict pilot regime is the next artifact-hardening
step (see Limitations in the main paper).

\section{Auxiliary Signal: Library-Health Contrasts under a Sonnet Pilot}
\label{app:lh-aux}

A separate preserved-source Sonnet pilot (four seeds, nine sessions)
applies the same harness, manifests, and verifier set as the
verified-subset Haiku pilot in the main paper. Per-tool source
preservation makes the capability-aligned conformance layer
computable, so library health (LH) -- combining correct reuse,
redundancy, utilisation, composition, regression, and the
hidden-test quality gate -- can be measured at the artifact level.
This is the same audit schema as the main pilot; only the model
snapshot and the source-preservation setting differ, which makes
the contrasts here a same-schema-different-model probe of the
framework's discriminative power, not a TC claim against the
verified-subset pilot.

\paragraph{Internal manifest labels.} The released \texttt{statistical\_report.json}
uses the project's internal names; the table below uses the public
naming from the main paper. The mapping is: \texttt{arise} $\rightarrow$
``Code-Evol'' (an internal preserved-source pilot lineage);
\texttt{evoskill} $\rightarrow$ ``Strategy-Only'' in the LH analysis
because LH does not depend on whether the strategy was executed;
\texttt{creator-style} and \texttt{no-evolution} keep the same labels.

\begin{table}[h]
\centering\small
\caption{Library-health contrasts in the Sonnet pilot
(Sonnet, $n{=}4$ seeds $\times$ 9 sessions; per-tool source preserved).
Paired hierarchical bootstrap over (seed, session), $10\,000$ resamples;
BH-adjusted bootstrap $p$. Generated by \texttt{scripts/stats\_report.py}.}
\label{tab:lh_stats_supp}
\begin{tabular}{lccc}
\toprule
\textbf{Contrast (LH)} & \textbf{$\Delta$} & \textbf{95\% CI} & \textbf{$p_{\mathrm{adj}}$} \\
\midrule
Code-Evol $-$ No-Evol           & 12.0$^{*}$ & [6.8, 17.0]  & 0.000 \\
CREATOR-style $-$ No-Evol       & 15.2$^{*}$ & [8.0, 23.4]  & 0.000 \\
Code-Evol $-$ Strategy-Only     & 10.4$^{*}$ & [6.4, 14.7]  & 0.000 \\
CREATOR-style $-$ Strategy-Only & 13.6$^{*}$ & [7.5, 21.3]  & 0.000 \\
\bottomrule
\end{tabular}
\\[2pt]\footnotesize $^{*}$95\% CI excludes 0. LH in percentage points.
\end{table}

These contrasts show that the framework can distinguish
tool-creating systems from non-creating systems on a measurable
artifact dimension (LH) when the underlying tool source is preserved
and per-tool conformance is aligned. Achieving the same alignment
under the strict Haiku pilot is the next artifact-hardening step. The
absolute LH levels and the magnitudes of these contrasts are
provider-specific and should not be transferred to the Haiku strict
pilot.

\onecolumn
\appendix
\section{Artifact Gallery}
\label{app:artifact_gallery}

The benchmark preserves every synthesised tool to disk under
\texttt{results\_full/<config>/tools/<session>/<name>.py} (plus a
\texttt{.meta.json} with the per-tool TQS scores). This appendix presents
a side-by-side gallery of representative tools from the two
tool-creating systems with multi-seed preserved source --- Code-Evol
(semantic-$Q$ run) and CREATOR-style --- to make the per-tool failure
modes the benchmark surfaces concrete. Per-dimension scores are shown
inline ($C$, $R$, $G$, $Q$ and the composite TQS); bold indicates
$\geq 0.5$.

\subsection{Code-Evol/Sonnet (semantic $Q$)}

15 tools were synthesised and preserved on disk across 9 sessions.
Below we show 5 representative tools spanning the quality
spectrum: correct exemplars, fragile high-$Q$/low-$C$ tools, duplicates,
and trivial (low-TQS) cases.

\noindent \textbf{Correct exemplar} \textit{(passes hidden tests with $C{\geq}0.5$)}\\
    \texttt{\textbf{http\_get\_with\_headers}} \footnotesize\textit{(api\_orchestration\_s1)}\quad $C{=}\textbf{0.50}$\;$R{=}\textbf{0.50}$\;$G{=}\textbf{0.50}$\;$Q{=}\textbf{1.00}$\;$\text{TQS}{=}\textbf{0.62}$
    \begin{lstlisting}[basicstyle=\scriptsize\ttfamily,language=Python,frame=none,xleftmargin=1em,breaklines=true,postbreak=\mbox{\textcolor{gray}{$\hookrightarrow$}\space},showstringspaces=false]
def http_get_with_headers(url: str, headers: dict) -> str:
    """Make an HTTP GET request with custom headers."""
    import urllib.request
    import urllib.error

    try:
        # Create request object with custom headers
        request = urllib.request.Request(url)
    # ...
    except Exception as e:
        return f"Error: {str(e)}"
\end{lstlisting}
    \vspace{0.8em}

\noindent \textbf{Correct exemplar} \textit{(passes hidden tests with $C{\geq}0.5$)}\\
    \texttt{\textbf{compute\_stats}} \footnotesize\textit{(numerical\_s1)}\quad $C{=}\textbf{0.50}$\;$R{=}\textbf{0.50}$\;$G{=}\textbf{0.50}$\;$Q{=}\textbf{1.00}$\;$\text{TQS}{=}\textbf{0.62}$
    \begin{lstlisting}[basicstyle=\scriptsize\ttfamily,language=Python,frame=none,xleftmargin=1em,breaklines=true,postbreak=\mbox{\textcolor{gray}{$\hookrightarrow$}\space},showstringspaces=false]
def compute_stats(data: list[float]) -> dict[str, float]:
    """Calculate basic statistical measures for a numerical dataset."""
    import math

    try:
        if not data:
            return {"error": "Empty dataset provided"}

    # ...
    except Exception as e:
        return {"error": f"Computation failed: {str(e)}"}
\end{lstlisting}
    \vspace{0.8em}

\noindent \textbf{Fragile} \textit{(high $Q$, low $C$ -- looks well-formed, fails on hidden inputs)}\\
    \texttt{\textbf{hmac\_sha256}} \footnotesize\textit{(api\_orchestration\_s1)}\quad $C{=}0.00$\;$R{=}\textbf{0.50}$\;$G{=}0.00$\;$Q{=}\textbf{1.00}$\;$\text{TQS}{=}0.38$
    \begin{lstlisting}[basicstyle=\scriptsize\ttfamily,language=Python,frame=none,xleftmargin=1em,breaklines=true,postbreak=\mbox{\textcolor{gray}{$\hookrightarrow$}\space},showstringspaces=false]
def hmac_sha256(secret: str, message: str) -> str:
    """Generate HMAC-SHA256 hash for authentication and security purposes."""
    import hmac
    import hashlib

    try:
        # Handle None inputs explicitly
        if secret is None or message is None:
    # ...
    except Exception as e:
        return f"Error generating HMAC-SHA256: {str(e)}"
\end{lstlisting}
    \vspace{0.8em}

\noindent \textbf{Fragile} \textit{(high $Q$, low $C$ -- looks well-formed, fails on hidden inputs)}\\
    \texttt{\textbf{compute\_crc32}} \footnotesize\textit{(data\_transform\_s5)}\quad $C{=}0.00$\;$R{=}\textbf{1.00}$\;$G{=}\textbf{1.00}$\;$Q{=}\textbf{1.00}$\;$\text{TQS}{=}\textbf{0.75}$
    \begin{lstlisting}[basicstyle=\scriptsize\ttfamily,language=Python,frame=none,xleftmargin=1em,breaklines=true,postbreak=\mbox{\textcolor{gray}{$\hookrightarrow$}\space},showstringspaces=false]
def compute_crc32(data: bytes) -> int:
    """Compute CRC32 checksum for data integrity verification."""
    import struct

    try:
        if not isinstance(data, bytes):
            return -1

    # ...
    except Exception:
        return -1
\end{lstlisting}
    \vspace{0.8em}

\noindent \textbf{Duplicate} \textit{(synthesised across multiple sessions)}\\
    \texttt{\textbf{decode\_base64}} \footnotesize\textit{(data\_transform\_s5)}\quad $C{=}\textbf{0.50}$\;$R{=}\textbf{0.50}$\;$G{=}\textbf{0.50}$\;$Q{=}\textbf{1.00}$\;$\text{TQS}{=}\textbf{0.62}$
    \begin{lstlisting}[basicstyle=\scriptsize\ttfamily,language=Python,frame=none,xleftmargin=1em,breaklines=true,postbreak=\mbox{\textcolor{gray}{$\hookrightarrow$}\space},showstringspaces=false]
def decode_base64(encoded_string: str) -> bytes:
    """Decode base64 encoded binary data into bytes."""
    import base64

    # Validate input type
    if not isinstance(encoded_string, str):
        raise TypeError(f"Expected string, got {type(encoded_string).__name__}")

    # ...
        # Return empty bytes on any decoding error (invalid base64, etc.)
        return b''
\end{lstlisting}
    \vspace{0.8em}

\subsection{CREATOR-style/Sonnet}

18 tools were synthesised and preserved on disk across 9 sessions.
Below we show 5 representative tools spanning the quality
spectrum: correct exemplars, fragile high-$Q$/low-$C$ tools, duplicates,
and trivial (low-TQS) cases.

\noindent \textbf{Correct exemplar} \textit{(passes hidden tests with $C{\geq}0.5$)}\\
    \texttt{\textbf{decode\_guardian\_minimal\_data}} \footnotesize\textit{(data\_transform\_s5)}\quad $C{=}\textbf{0.50}$\;$R{=}\textbf{0.50}$\;$G{=}\textbf{0.50}$\;$Q{=}\textbf{0.89}$\;$\text{TQS}{=}\textbf{0.60}$
    \begin{lstlisting}[basicstyle=\scriptsize\ttfamily,language=Python,frame=none,xleftmargin=1em,breaklines=true,postbreak=\mbox{\textcolor{gray}{$\hookrightarrow$}\space},showstringspaces=false]
def decode_guardian_minimal_data(base64_data):
    """Decode minimal GUARDIAN data format handling single blocks and edge cases."""
    import base64
    import struct

    try:
        # Decode base64 data
        raw_data = base64.b64decode(base64_data)
    # ...
    except Exception as e:
        return {'text': '', 'integrity': False, 'blocks_processed': 0, 'details': f'Error: {str(e)}'}
\end{lstlisting}
    \vspace{0.8em}

\noindent \textbf{Correct exemplar} \textit{(passes hidden tests with $C{\geq}0.5$)}\\
    \texttt{\textbf{decode\_qlog\_with\_continuations}} \footnotesize\textit{(data\_transform\_s3)}\quad $C{=}\textbf{0.50}$\;$R{=}\textbf{0.50}$\;$G{=}\textbf{0.50}$\;$Q{=}\textbf{0.89}$\;$\text{TQS}{=}\textbf{0.60}$
    \begin{lstlisting}[basicstyle=\scriptsize\ttfamily,language=Python,frame=none,xleftmargin=1em,breaklines=true,postbreak=\mbox{\textcolor{gray}{$\hookrightarrow$}\space},showstringspaces=false]
def decode_qlog_with_continuations(base64_data):
    """Decode QLOG data and merge continuation entries with their parent entries."""
    import base64
    import struct
    from datetime import datetime, timezone

    # Decode base64 data
    binary_data = base64.b64decode(base64_data)
    # ...

    return merged_entries
\end{lstlisting}
    \vspace{0.8em}

\noindent \textbf{Fragile} \textit{(high $Q$, low $C$ -- looks well-formed, fails on hidden inputs)}\\
    \texttt{\textbf{fetch\_all\_users\_with\_pagination}} \footnotesize\textit{(api\_orchestration\_s1)}\quad $C{=}0.00$\;$R{=}\textbf{0.50}$\;$G{=}\textbf{1.00}$\;$Q{=}\textbf{0.89}$\;$\text{TQS}{=}\textbf{0.60}$
    \begin{lstlisting}[basicstyle=\scriptsize\ttfamily,language=Python,frame=none,xleftmargin=1em,breaklines=true,postbreak=\mbox{\textcolor{gray}{$\hookrightarrow$}\space},showstringspaces=false]
def fetch_all_users_with_pagination(base_url="http://127.0.0.1:18080/api/users"):
    """Utility: Fetches all users from a paginated API endpoint using cursor-based pagination.
Continues fetching until has_more is False, collecting all user data across pages."""
    import urllib.request
    import urllib.parse
    import json

    all_users = []
    cursor = None
    # ...
        'names': names
    }
\end{lstlisting}
    \vspace{0.8em}

\noindent \textbf{Fragile} \textit{(high $Q$, low $C$ -- looks well-formed, fails on hidden inputs)}\\
    \texttt{\textbf{decode\_qlog\_binary\_data}} \footnotesize\textit{(data\_transform\_s3)}\quad $C{=}0.00$\;$R{=}\textbf{0.50}$\;$G{=}0.00$\;$Q{=}\textbf{0.78}$\;$\text{TQS}{=}0.32$
    \begin{lstlisting}[basicstyle=\scriptsize\ttfamily,language=Python,frame=none,xleftmargin=1em,breaklines=true,postbreak=\mbox{\textcolor{gray}{$\hookrightarrow$}\space},showstringspaces=false]
def decode_qlog_binary_data(base64_data):
    """Decode QLOG (Quantized Log Format) binary data into structured log records."""
    import base64
    import struct
    from datetime import datetime, timezone

    # Decode base64 data
    binary_data = base64.b64decode(base64_data)
    # ...

    return records
\end{lstlisting}
    \vspace{0.8em}

\noindent \textbf{Trivial} \textit{(low overall TQS)}\\
    \texttt{\textbf{deserialize\_tpack}} \footnotesize\textit{(data\_transform\_s4)}\quad $C{=}0.00$\;$R{=}\textbf{0.50}$\;$G{=}0.00$\;$Q{=}\textbf{0.67}$\;$\text{TQS}{=}0.29$
    \begin{lstlisting}[basicstyle=\scriptsize\ttfamily,language=Python,frame=none,xleftmargin=1em,breaklines=true,postbreak=\mbox{\textcolor{gray}{$\hookrightarrow$}\space},showstringspaces=false]
def deserialize_tpack(base64_data):
    """Deserialize TPACK (Tagged Pack Format) binary data into Python objects."""
    import base64
    import struct

    def parse_varint(data, offset):
        result = 0
        shift = 0
    # ...
    result, _ = parse_value(binary_data, 0)
    return result
\end{lstlisting}
    \vspace{0.8em}


\end{document}
